\documentclass[11pt]{article}
\usepackage[T1]{fontenc}
\usepackage[utf8]{inputenc}
\usepackage{amsmath,amssymb,amsthm}
\usepackage[margin=2.5cm]{geometry}
\usepackage{booktabs}
\setlength{\parindent}{1.4em}
\setlength{\parskip}{0pt}

\newtheorem{claim}{Claim}
\newtheorem{remark}{Remark}

\begin{document}

% =============== PAGE 1 ===============
\noindent{\LARGE\textbf{Architecture {\large $>$} Population:}}\\[0.15em]
\noindent{\Large\textbf{A Research Program for the Inter-Agent Mesoscale}}\\[0.4em]
\noindent{\small Hakvin Vosteen \hfill 2026 \qquad \textit{Draft, v0.1.}}
\vspace{0.6cm}

\noindent\textbf{Abstract.}\quad The thirteen working pieces in this corpus
repeat a single methodological move: a phenomenon ordinarily read as moral,
cultural, or characterological is rewritten as the equilibrium of a mechanism
with a threshold. We state the move as a research program organised around
three claims. First, what looks like a property of \textit{people} is usually
a property of a \textit{mechanism}; the recurrent objects are operators,
sufficient statistics, and threshold games. Second, architecture beats
population: replace the users while holding the mechanism fixed and the
equilibrium is preserved; replace the mechanism and the equilibrium changes.
Third, the scale at which most of contemporary subjective experience is
decided --- the inter-agent mesoscale --- is methodologically nearly empty
compared to the scales below the human (atoms, cells, neurons) and above
(climate, markets-as-aggregate). The marginal return on formalising one
previously unmodelled inter-agent mechanism is therefore large. We address
the standard reflexive objection (modelling perturbs the modelled system) by
distinguishing three modes of impossibility --- information-theoretic,
reflexive, and conceptual --- and showing, with anchors in the corpus, that
each admits a different attack.

\bigskip\noindent\textbf{1.\quad The move.}\quad Each piece in this corpus
takes a folk-psychological reading of a phenomenon and replaces it with a
mechanism. \textit{Burnout Is Not Too Much Work} replaces a moral reading
(weakness, lack of grit) with the bifurcation of a friction-output ratio at
$\rho^* = 1$. \textit{Ghosting Isn't Cowardice} replaces a generational or
character reading with the subgame-perfect equilibrium of a sequential exit
game indexed by a silence-cost ratio $\sigma_R$. \textit{ADHD Is Not a
Deficit} replaces a deficit reading with a different attractor in a phase
space whose stationary distribution is bimodal under a free-energy
generative model. \textit{Personality Is Not Who You Are} replaces a vector
of traits with a four-dimensional cost tensor under a context operator.
\textit{Cartel or Creative Destruction?} replaces concentration-index
conflation with eigenvector asymmetry. The move is the same in each case:
take what reads as a property of the agent, and locate it in the structure
the agent inhabits.

\medskip\noindent Two claims sit underneath the move.

\begin{claim}[Mechanism, not person]
Phenomena that read as moral, cultural, or characterological are typically
equilibria of mechanisms with thresholds. The dominant variable is not the
agent's interior; it is the cost structure, signal architecture, or
selection rule under which the agent acts.
\end{claim}

\begin{claim}[Architecture, not population]
Hold the mechanism fixed and replace the population: the equilibrium is
preserved. Hold the population fixed and replace the mechanism: the
equilibrium changes. Behaviour traceable to the population alone is a small
residual on top of architecture-driven dynamics.
\end{claim}

\noindent Claim~1 is identification of the dominant explanatory variable.
Claim~2 is the operational counterfactual that licenses Claim~1: the
population is held fixed and the architecture is varied. \textit{Ghosting
Isn't Cowardice} states the counterfactual explicitly --- replacing the user
pool with people from 1995 under the same dating-app architecture produces
the same equilibrium --- and the move recurs across the corpus.

\newpage

% =============== PAGE 2 ===============
\noindent\textbf{2.\quad Why this scale.}\quad Classical natural science has
measured the scales below the human (atoms, cells, neurons) and above the
human (climate, markets-as-aggregate, populations) to high resolution. The
scale \textit{between} two or $N$ agents interacting through an architecture
is, by comparison, methodologically nearly empty. There is no canonical
formal language for it. The most-cited models live in microeconomics (game
theory) and computer science (mechanism design) but are rarely applied
outside their disciplines, and the phenomena most affected by inter-agent
dynamics --- platform behaviour, dating, peer review, attention extraction,
identity signalling --- are typically left to qualitative discourse.

\medskip\noindent The marginal return on understanding one more electron is,
for the purposes of subjective experience, near zero. The marginal return on
formalising one previously unmodelled inter-agent mechanism is large,
because the prior is not ``well-established theory needs an additional
decimal place'' but ``no formal theory exists.'' This is the slot the corpus
operates in. It is also the structural reason that pieces of mathematics
typically considered elementary --- a multiplicative cascade, a free-energy
bound, an eigenvector asymmetry --- yield real explanatory traction in this
regime: the comparison is not against tighter models but against no model at
all.

\medskip\noindent\textit{A note on what is and is not new.} The argument that
inter-agent structure carries explanatory weight is not new in any single
discipline; it is the operating assumption of mechanism design, of much of
sociology, of evolutionary game theory. What is being claimed is narrower.
First, that across these disciplines the move is consistently localised to
its native vocabulary, so the same structural lemma --- a threshold game,
say --- gets re-derived in three places without recognition. Second, that
many of the phenomena that fill ordinary life (ghosting, burnout, identity
broadcast on platforms, the felt arbitrariness of recommendation) sit
between disciplines and are therefore picked up by none. The corpus reads
across these boundaries.

\bigskip\noindent\textbf{3.\quad Modelling under observer effects.}\quad The
standard objection to formalising inter-agent mechanisms is reflexive:
publishing a model of $\sigma_R$ changes how people ghost; publishing the
platform food chain changes how creators read their own role; the act of
measurement is itself the perturbation. The objection is real. It is not
grounds to refuse the work, because \textit{impossibility} is not a single
condition --- it has facets, and each admits a different attack.

\begin{remark}[Three modes of impossibility]
\hfill\\[0.3em]
\renewcommand{\arraystretch}{1.25}
\begin{tabular}{@{}p{0.18\linewidth}p{0.36\linewidth}p{0.40\linewidth}@{}}
\toprule
\textbf{Mode} & \textbf{Difficulty} & \textbf{Attack} \\
\midrule
Information-theoretic
  & Capturing the full state would require infinite or impossibly clean data
  & Reduce to a sufficient statistic; the relevant scalar can be finite even
    when the underlying process is high-dimensional \\
Reflexive
  & The model perturbs the system being modelled
  & Publish the model jointly with the architectural change that defines
    the new equilibrium, so the perturbation is directed rather than
    incidental \\
Conceptual
  & No agreed-upon target object exists
  & Show that the missing scalar ground truth was never the right object;
    supply the structured object that does the work \\
\bottomrule
\end{tabular}
\end{remark}

\medskip\noindent\textit{Information-theoretic.}\quad \textit{Cartel or
Creative Destruction?} faces precisely this difficulty: standard
concentration indices are degenerate across two welfare-opposite regimes
(collusive cartel and Schumpeterian innovation-driven monopoly). The attack
is to identify a sufficient statistic --- the eigenvector asymmetry of the
residual price covariance --- that separates the regimes with
Cohen's $d > 70$. The scalar is finite even though the underlying market
state is not.

\medskip\noindent\textit{Reflexive.}\quad \textit{Ghosting Isn't Cowardice}
models $\sigma_R$ and immediately proposes the architectural lever (e.g.\
Hinge's ``Your Turn'' feature) that raises $c_{\text{silence}}$ and flips
the equilibrium. The model is not a passive observation but a directed
perturbation: it ships with the change that justifies it. \textit{The
TikTok Peer Review} does the analogous move at second order, modelling the
mechanism that produces models. In both cases reflexivity is acknowledged
and converted into design.

\newpage

% =============== PAGE 3 ===============
\noindent\textit{Conceptual.}\quad \textit{Personality Is Not Who You Are.
It Is Your Cost Tensor} is the cleanest case. The Big Five framework
presupposes a five-dimensional scalar object whose component values
describe a person stably across context. The cost-tensor model shows this
object cannot exist as specified, because the relevant quantity is
context-dependent in a way no projection $\Pi_5$ can preserve. The apparent
unmeasurability of personality is a misspecification of the target object,
not a fundamental limit of measurement. Once the object is replaced by a
cost tensor under a context operator $L_W$, it is identifiable; the
intervention mechanism (operator sequences with eigenvalues less than one)
is also identified.

\bigskip\noindent\textbf{4.\quad What this program is and is not.}\quad The
program is not a claim that subjective experience reduces to mechanism. The
cost-tensor paper itself acknowledges $\Lambda_{\text{trauma}}$ as a
strictly individual operator with at least one eigenvalue $\gg 1$; some
variance is genuinely person-shaped. The claim is narrower: that the
variance commonly attributed to character is largely architectural, and
that the architectural component is the under-modelled component.
``Architecture $>$ population'' is a statement about marginal explanatory
power, not about ontology.

\medskip\noindent The program is also not a claim that all formalisation is
equally tractable. Cartel detection admits a clean sufficient statistic;
subjective wellbeing under social comparison probably does not. The three
modes of impossibility are not a guarantee that every reflexive system can
be modelled, only that the standard rejection (``you will perturb it,
therefore do not'') over-generalises a real local difficulty into a
fictitious global one.

\medskip\noindent What unites the corpus, then, is not a domain (which
ranges from drug discovery to dating apps) and not a single mathematical
object (operators, multiplicative cascades, free-energy bounds, threshold
games all appear), but a methodological commitment: locate the dominant
variable in the architecture, not the population; treat reflexive
perturbation as something to be designed rather than avoided; and, where
the standard target object is conceptually misspecified, supply the right
one rather than concede unmeasurability.

\bigskip\noindent\textbf{5.\quad Status.}\quad This is a draft (v0.1)
written in conversation with the corpus, not in front of it. It will be
revised as the corpus grows and as the connections between pieces become
load-bearing. The next obvious extensions are (i) a structural lemma
explicitly unifying the threshold-bifurcation skeleton across
\textit{Burnout}, \textit{Ghosting}, \textit{ADHD Attractor}, and the
\textit{Cost Tensor}; (ii) a precise statement of the
``architecture-not-population'' identification result, with its scope and
failure modes; (iii) a fourth mode of impossibility, computational, for
systems where a tractable model exists in principle but not in practice.

\bigskip\noindent\textbf{Cited works in this corpus.}\quad
\textit{Burnout Is Not Too Much Work, v2} (2026-05-06).\quad
\textit{Ghosting Isn't Cowardice. It's Math.} (2026-04-24).\quad
\textit{ADHD Is Not a Deficit. It's a Different Attractor.} (2026-04-24).\quad
\textit{Personality Is Not Who You Are. It Is Your Cost Tensor.}
(2026-05-06).\quad
\textit{Cartel or Creative Destruction? Eigenvector Asymmetry as a Spectral
Diagnostic for Schumpeterian Market Structure} (2026-04-22).\quad
\textit{The TikTok Peer Review} (2026-04-26).\quad
\textit{The Platform Has a Food Chain. You Are In It.} (2026-05-06).\quad
All available at \texttt{hakvinv.github.io/h.vosteen-research-/}.

\end{document}
